\documentclass[11pt,a4paper]{article}
\usepackage[utf8]{inputenc}
\usepackage[spanish]{babel}
\usepackage[T1]{fontenc}
\usepackage{geometry}
\geometry{margin=2cm}
\usepackage{xcolor}
\usepackage{listings}
\usepackage{tcolorbox}
\tcbuselibrary{most}
\usepackage{booktabs}
\usepackage{array}
\usepackage{hyperref}
\usepackage{fancyhdr}
\usepackage{titlesec}
\usepackage{multicol}
\usepackage{amssymb}
\usepackage{pifont}

% Símbolos D&D sin fontawesome
\newcommand{\iconCofre}{$\bigstar$}
\newcommand{\iconTrampa}{$\danger$}
\newcommand{\iconHechizo}{$\lozenge$}
\newcommand{\iconPergamino}{$\triangleright$}
\newcommand{\iconMision}{$\blacksquare$}
\newcommand{\iconRecompensa}{$\checkmark$}
\newcommand{\iconLore}{$\S$}
\newcommand{\danger}{\text{\ding{56}}}

% ─── COLORES ──────────────────────────────────────────────────────────────────
\definecolor{vuegreen}{RGB}{66,184,131}
\definecolor{vuedark}{RGB}{35,120,85}
\definecolor{golddnd}{RGB}{193,145,32}
\definecolor{redrune}{RGB}{180,40,40}
\definecolor{purplemage}{RGB}{100,60,160}
\definecolor{codebg}{RGB}{245,245,245}
\definecolor{alertred}{RGB}{220,50,50}
\definecolor{tipgreen}{RGB}{39,174,96}
\definecolor{warnambel}{RGB}{230,126,34}
\definecolor{darkgray}{RGB}{50,50,50}
\definecolor{tsblue}{RGB}{0,122,204}
\definecolor{parchment}{RGB}{255,248,220}
\definecolor{parchmentborder}{RGB}{160,120,60}

% ─── LISTING: TypeScript/Vue ──────────────────────────────────────────────────
\lstdefinestyle{vue}{
  language=Java,
  backgroundcolor=\color{codebg},
  basicstyle=\ttfamily\small,
  keywordstyle=\color{vuegreen}\bfseries,
  stringstyle=\color{teal},
  commentstyle=\color{gray}\itshape,
  numbers=left, numberstyle=\tiny\color{gray},
  frame=single, framerule=0.5pt,
  rulecolor=\color{vuegreen!50},
  breaklines=true, breakatwhitespace=true,
  tabsize=2, showstringspaces=false,
  morekeywords={ref,reactive,computed,watch,watchEffect,
    onMounted,onUnmounted,onBeforeMount,onUpdated,
    defineProps,defineEmits,defineExpose,
    provide,inject,useSlots,useAttrs,
    readonly,toRef,toRefs,
    const,let,interface,type,async,await,
    string,number,boolean,void,null,
    import,from,export,default,
    storeToRefs,defineStore,acceptHMRUpdate}
}

\lstdefinestyle{bash}{
  language=bash,
  backgroundcolor=\color{codebg},
  basicstyle=\ttfamily\small,
  frame=single, framerule=0.4pt,
  rulecolor=\color{gray!40},
  breaklines=true, numbers=none
}

% ─── CAJAS TEMÁTICAS D&D ──────────────────────────────────────────────────────
% Cofre del tesoro = Tip
\newtcolorbox{cofre}{
  colback=golddnd!12, colframe=golddnd,
  title={\ding{72}\; \textbf{Cofre del Tesoro}},
  fonttitle=\small\bfseries
}

% Trampa = Alerta / Error común
\newtcolorbox{trampa}{
  colback=redrune!10, colframe=redrune,
  title={\ding{56}\; \textbf{¡Trampa!}},
  fonttitle=\small\bfseries
}

% Hechizo = Concepto clave
\newtcolorbox{hechizo}[1]{
  colback=purplemage!8, colframe=purplemage,
  title={\ding{168}\; \textbf{Hechizo: #1}},
  fonttitle=\small\bfseries
}

% Pergamino = Código / ejemplo
\newtcolorbox{pergamino}{
  colback=vuegreen!6, colframe=vuegreen,
  title={\ding{228}\; \textbf{Pergamino de código}},
  fonttitle=\small\bfseries
}

% Misión = Ejercicio
\newtcolorbox{mision}[1]{
  colback=warnambel!8, colframe=warnambel,
  title={\ding{110}\; \textbf{Misión #1}},
  fonttitle=\small\bfseries
}

% Recompensa = Solución
\newtcolorbox{recompensa}[1]{
  colback=tipgreen!8, colframe=tipgreen!70,
  title={\ding{52}\; \textbf{Recompensa #1}},
  fonttitle=\small\bfseries
}

% Lore = Contexto / nota de mundo
\newtcolorbox{lore}{
  colback=parchment, colframe=parchmentborder,
  title={\S\; \textbf{Lore del Mundo}},
  fonttitle=\small\bfseries\color{parchmentborder}
}

% ─── ENCABEZADO ───────────────────────────────────────────────────────────────
\pagestyle{fancy}
\fancyhf{}
\lhead{\textcolor{vuegreen}{\textbf{Vue.js — El Libro del Aventurero}}}
\rhead{\textcolor{gray}{Guía de Componentes y Reactividad}}
\cfoot{\thepage}
\setlength{\headheight}{15pt}

% ─── SECCIONES ────────────────────────────────────────────────────────────────
% Sin numeración automática: los títulos incluyen "Nivel N" manualmente
\setcounter{secnumdepth}{0}
\titleformat{\section}
  {\large\bfseries\color{vuegreen}}{}{0em}{}[\titlerule]
\titleformat{\subsection}
  {\normalsize\bfseries\color{darkgray}}{}{0em}{}

% ══════════════════════════════════════════════════════════════════════════════
\begin{document}

% ─── PORTADA ──────────────────────────────────────────────────────────────────
\begin{titlepage}
  \centering
  \vspace*{1.5cm}

  {\small\color{gray}\textit{— El Libro del Aventurero —}\par}
  \vspace{0.4cm}
  {\Huge\bfseries\color{vuegreen} Vue.js\par}
  {\large\bfseries\color{vuedark} Guía del Héroe Frontend\par}
  \vspace{0.3cm}
  \textcolor{golddnd}{\rule{10cm}{2pt}}
  \vspace{0.8cm}

  \begin{tcolorbox}[
    colback=parchment, colframe=parchmentborder,
    width=11cm
  ]
    \centering
    \textit{``En las tierras de Vue.js, todo componente es un héroe.
    Aprenderá reactividad, dominará los composables y enfrentará
    el desafío final: construir una UI que cobros recurrentes
    ni el banco más temible logrará romper.''}\par
    \vspace{6pt}
    {\small — El Libro del Aventurero, Prólogo}
  \end{tcolorbox}

  \vspace{0.8cm}
  \begin{tcolorbox}[
    colback=vuegreen!5, colframe=vuegreen, width=10cm
  ]
    \centering
    \textbf{Stack objetivo:} \texttt{Vue 3 + TypeScript + Pinia}\\[4pt]
    \textbf{Versión:} Vue 3.4+ (Composition API)\\[4pt]
    \textbf{Compatible con:} Stack Toku (Vue 3 + FastAPI)
  \end{tcolorbox}

  \vspace{1cm}
  {\large Árbol de Habilidades\par}
  \vspace{0.3cm}
  \begin{tabular}{ll}
    \textcolor{vuegreen}{\ding{108}} &
      Nivel 1 — Reactividad: \texttt{ref}, \texttt{reactive}, \texttt{computed}\\
    \textcolor{vuegreen}{\ding{108}} &
      Nivel 2 — Ciclo de vida y \texttt{watch}\\
    \textcolor{vuegreen}{\ding{108}} &
      Nivel 3 — Componentes, Props y Emits\\
    \textcolor{vuegreen}{\ding{108}} &
      Nivel 4 — Composables: pergaminos reutilizables\\
    \textcolor{golddnd}{\ding{72}}           &
      Nivel 5 — Pinia: el Grimorio del Estado Global\\
    \textcolor{golddnd}{\ding{72}}           &
      Nivel 6 — Patrones avanzados y Slots\\
    \textcolor{redrune}{\ding{58}}         &
      Nivel 7 — Misiones finales (ejercicios integrados)\\
  \end{tabular}

  \vfill
  {\small\color{gray} Marzo 2026 — Estudio Personal\par}
\end{titlepage}

\tableofcontents
\newpage

% ══════════════════════════════════════════════════════════════════════════════
% NIVEL 1: REACTIVIDAD
% ══════════════════════════════════════════════════════════════════════════════
\section{Nivel 1 — Reactividad: el Corazón del Héroe}

\begin{lore}
En las antiguas tierras de jQuery, el DOM se manipulaba a mano, fila por fila.
Vue.js cambió las reglas: declara \textit{qué} debe verse, y el motor reactivo
se encarga de \textit{cómo} actualizarlo. No hay varitas mágicas --- hay
\texttt{ref()} y \texttt{reactive()}.
\end{lore}

\subsection{ref() — La Gema de Memoria}

\begin{hechizo}{ref()}
Envuelve un valor primitivo (número, string, boolean) en un objeto reactivo.
Dentro del \texttt{<script>} se accede con \texttt{.value}.
En el \texttt{<template>} Vue lo desenvuelve automáticamente.
\end{hechizo}

\begin{pergamino}
\begin{lstlisting}[style=vue]
<script setup lang="ts">
import { ref } from 'vue'

// Tipos simples
const gold = ref<number>(0)
const heroName = ref<string>('Fabian')
const isAlive = ref<boolean>(true)

// Tipos complejos (array, objeto)
const inventory = ref<string[]>([])
const stats = ref<{ hp: number; mp: number }>({
  hp: 100,
  mp: 50
})

// Modificar: siempre con .value en <script>
function collectGold(amount: number) {
  gold.value += amount
  inventory.value.push(`${amount} oro`)
}

// Leer el valor
console.log(gold.value) // 0
</script>

<template>
  <!-- En el template: sin .value -->
  <p>{{ heroName }} tiene {{ gold }} de oro</p>
  <button @click="collectGold(10)">Recolectar</button>
  <ul>
    <li v-for="item in inventory" :key="item">
      {{ item }}
    </li>
  </ul>
</template>
\end{lstlisting}
\end{pergamino}

\begin{trampa}
\textbf{Error frecuente:} olvidar \texttt{.value} en el \texttt{<script>}.
\begin{lstlisting}[style=vue,numbers=none]
// MAL: asigna directamente (rompe la reactividad)
gold = 50

// BIEN: modifica el valor reactivo
gold.value = 50

// MAL: reeemplaza el array (pierde la referencia)
inventory.value = [...inventory.value, 'espada']
// BIEN para arrays:
inventory.value.push('espada')
\end{lstlisting}
\end{trampa}

\subsection{reactive() — El Objeto Vivo}

\begin{hechizo}{reactive()}
Hace reactivo un objeto completo. No necesita \texttt{.value} para
acceder a sus propiedades. Ideal para agrupar estado relacionado.
\end{hechizo}

\begin{pergamino}
\begin{lstlisting}[style=vue]
<script setup lang="ts">
import { reactive, toRefs } from 'vue'

interface HeroState {
  name: string
  hp: number
  mp: number
  inventory: string[]
  isLoading: boolean
  error: string | null
}

const state = reactive<HeroState>({
  name: 'Fabian',
  hp: 100,
  mp: 50,
  inventory: [],
  isLoading: false,
  error: null
})

// Acceso directo, sin .value
function takeDamage(amount: number) {
  state.hp -= amount
  if (state.hp <= 0) {
    state.error = 'El heroe ha caido'
  }
}

// toRefs: destruir reactive sin perder reactividad
const { name, hp, mp } = toRefs(state)
// Ahora name.value, hp.value funcionan como ref
</script>
\end{lstlisting}
\end{pergamino}

\begin{cofre}
\textbf{ref vs reactive:} usa \texttt{ref()} para valores simples y
cuando necesitas pasar el valor como parámetro (mantiene la referencia).
Usa \texttt{reactive()} para agrupar estado relacionado de un componente.
En la práctica, muchos equipos usan \texttt{ref()} para todo para mayor
consistencia.
\end{cofre}

\subsection{computed() — El Oráculo}

\begin{hechizo}{computed()}
Deriva un valor de otros valores reactivos. Se recalcula
automáticamente solo cuando cambia alguna dependencia. Nunca
modifica --- solo lee.
\end{hechizo}

\begin{pergamino}
\begin{lstlisting}[style=vue]
<script setup lang="ts">
import { ref, computed } from 'vue'

const price = ref<number>(10000)  // en centavos
const quantity = ref<number>(3)
const discountPct = ref<number>(10) // 10%

// computed: se recalcula cuando price o quantity cambian
const subtotal = computed(() =>
  price.value * quantity.value
)

// computed con dependencias multiples
const discount = computed(() =>
  subtotal.value * (discountPct.value / 100)
)

const total = computed(() =>
  subtotal.value - discount.value
)

// computed writable (get + set)
const displayPrice = computed({
  get: () => price.value / 100,        // centavos -> pesos
  set: (pesos: number) => {
    price.value = Math.round(pesos * 100) // pesos -> centavos
  }
})
</script>

<template>
  <p>Subtotal: ${{ (subtotal / 100).toFixed(2) }}</p>
  <p>Descuento: -${{ (discount / 100).toFixed(2) }}</p>
  <p><strong>Total: ${{ (total / 100).toFixed(2) }}</strong></p>
</template>
\end{lstlisting}
\end{pergamino}

\begin{trampa}
No hagas efectos secundarios dentro de \texttt{computed} (llamadas a API,
mutaciones). Para eso existe \texttt{watch}. Un computed que modifica
estado es una trampa garantizada de loops infinitos.
\end{trampa}

% ══════════════════════════════════════════════════════════════════════════════
% NIVEL 2: CICLO DE VIDA Y WATCH
% ══════════════════════════════════════════════════════════════════════════════
\section{Nivel 2 — Ciclo de Vida y Centinelas}

\subsection{Ciclo de vida de un componente}

\begin{center}
\begin{tabular}{p{3.5cm} p{9cm}}
\toprule
\textbf{Hook} & \textbf{Cuándo se ejecuta} \\
\midrule
\texttt{onBeforeMount} & Antes de que el DOM exista. No acceder a elementos. \\
\texttt{onMounted}     & DOM listo. Ideal para llamadas a API, subscripciones. \\
\texttt{onBeforeUpdate}& Justo antes de re-render por cambio reactivo. \\
\texttt{onUpdated}     & Después del re-render. Raro usarlo. \\
\texttt{onBeforeUnmount}& Componente a punto de destruirse. Cancelar suscripciones. \\
\texttt{onUnmounted}   & Componente destruido. Limpiar timers, listeners. \\
\texttt{onErrorCaptured}& Error en componente hijo. Para error boundaries. \\
\bottomrule
\end{tabular}
\end{center}

\begin{pergamino}
\begin{lstlisting}[style=vue]
<script setup lang="ts">
import { ref, onMounted, onUnmounted } from 'vue'

const charges = ref<Charge[]>([])
const isLoading = ref(false)
const error = ref<string | null>(null)
let pollingInterval: ReturnType<typeof setInterval> | null = null

// onMounted: el momento de actuar
onMounted(async () => {
  await loadCharges()

  // Polling cada 30 segundos
  pollingInterval = setInterval(loadCharges, 30_000)
})

// onUnmounted: limpiar lo que levantamos
onUnmounted(() => {
  if (pollingInterval) {
    clearInterval(pollingInterval)
  }
})

async function loadCharges() {
  isLoading.value = true
  error.value = null
  try {
    const resp = await fetch('/api/charges')
    charges.value = await resp.json()
  } catch {
    error.value = 'Error al cargar cobros'
  } finally {
    isLoading.value = false
  }
}
</script>
\end{lstlisting}
\end{pergamino}

\subsection{watch() y watchEffect() — Los Centinelas}

\begin{pergamino}
\begin{lstlisting}[style=vue]
<script setup lang="ts">
import { ref, watch, watchEffect } from 'vue'

const searchQuery = ref('')
const page = ref(1)
const results = ref<any[]>([])

// watch: observa fuentes especificas
// Se ejecuta solo cuando searchQuery cambia
watch(searchQuery, async (newVal, oldVal) => {
  if (newVal === oldVal) return
  page.value = 1  // resetear pagina al buscar
  await fetchResults(newVal, 1)
}, { debounce: 300 })  // esperar 300ms

// watch con multiples fuentes
watch([searchQuery, page], async ([query, p]) => {
  await fetchResults(query, p)
})

// watchEffect: detecta dependencias automaticamente
// Se ejecuta inmediatamente y cada vez que cambia algo
watchEffect(async () => {
  // Vue detecta que usa searchQuery.value y page.value
  results.value = await fetchResults(
    searchQuery.value,
    page.value
  )
})

// watch con opciones
watch(searchQuery, handler, {
  immediate: true,  // ejecutar inmediatamente al montar
  deep: true,       // observar cambios profundos en objetos
  flush: 'post',    // despues del render
})
</script>
\end{lstlisting}
\end{pergamino}

\begin{cofre}
\textbf{¿watch o watchEffect?}
Usa \texttt{watch} cuando necesitas saber el valor anterior o cuando
quieres control explícito de qué observar.
Usa \texttt{watchEffect} cuando la lógica es corta y quieres que
Vue detecte automáticamente las dependencias.
\end{cofre}

% ══════════════════════════════════════════════════════════════════════════════
% NIVEL 3: COMPONENTES
% ══════════════════════════════════════════════════════════════════════════════
\section{Nivel 3 — Componentes: Héroes y Mensajeros}

\begin{lore}
Un componente bien diseñado es como un mercenario especializado: sabe
exactamente qué necesita (props), qué puede comunicar (emits) y no
guarda secretos que no le pertenecen (sin estado global innecesario).
\end{lore}

\subsection{Props — Los Mensajeros}

\begin{pergamino}
\begin{lstlisting}[style=vue]
<!-- ChargeCard.vue -->
<script setup lang="ts">
interface Charge {
  id: string
  amount: number
  status: 'pending' | 'succeeded' | 'failed'
  mandateId: string
  createdAt: string
}

// defineProps con TypeScript (sin runtime validators)
interface Props {
  charge: Charge
  showDetails?: boolean  // opcional con default
  currency?: 'CLP' | 'MXN' | 'BRL'
}

const props = withDefaults(defineProps<Props>(), {
  showDetails: false,
  currency: 'CLP'
})

// computed derivado de una prop
const formattedAmount = computed(() =>
  (props.charge.amount / 100).toLocaleString('es-CL', {
    style: 'currency',
    currency: props.currency
  })
)
</script>

<template>
  <div class="charge-card"
       :class="`status-${charge.status}`">
    <span>{{ charge.id }}</span>
    <span>{{ formattedAmount }}</span>
    <span>{{ charge.status }}</span>

    <!-- Mostrar detalles solo si se pide -->
    <div v-if="showDetails">
      <p>Mandato: {{ charge.mandateId }}</p>
      <p>Fecha: {{ charge.createdAt }}</p>
    </div>
  </div>
</template>
\end{lstlisting}
\end{pergamino}

\subsection{Emits — Las Señales de Humo}

\begin{pergamino}
\begin{lstlisting}[style=vue]
<!-- ChargeActions.vue -->
<script setup lang="ts">
interface Props {
  chargeId: string
  canRetry: boolean
}

const props = defineProps<Props>()

// Emits tipados: nombre del evento -> tipos de payload
const emit = defineEmits<{
  'retry-charge': [chargeId: string]
  'view-details': [chargeId: string]
  'export-pdf': []
}>()

async function handleRetry() {
  // validar antes de emitir
  if (!props.canRetry) return
  emit('retry-charge', props.chargeId)
}
</script>

<template>
  <div class="actions">
    <button
      :disabled="!canRetry"
      @click="handleRetry"
    >
      Reintentar cobro
    </button>
    <button @click="emit('view-details', chargeId)">
      Ver detalle
    </button>
    <button @click="emit('export-pdf')">
      Exportar PDF
    </button>
  </div>
</template>
\end{lstlisting}
\end{pergamino}

\begin{pergamino}
\begin{lstlisting}[style=vue]
<!-- Padre: ChargeList.vue -->
<script setup lang="ts">
import ChargeActions from './ChargeActions.vue'

async function onRetry(chargeId: string) {
  await retryCharge(chargeId)
  await refreshList()
}
</script>

<template>
  <!-- Escuchar emits con @nombre-evento -->
  <ChargeActions
    :charge-id="charge.id"
    :can-retry="charge.status === 'failed_retryable'"
    @retry-charge="onRetry"
    @view-details="router.push(`/charges/${$event}`)"
    @export-pdf="downloadPDF()"
  />
</template>
\end{lstlisting}
\end{pergamino}

\begin{trampa}
Las props son \textbf{de solo lectura}. Jamás las modifiques directamente
en el componente hijo --- Vue lanzará una advertencia. Si necesitas un
estado local basado en una prop, crea un \texttt{ref} derivado:
\begin{lstlisting}[style=vue,numbers=none]
const props = defineProps<{ initialValue: number }>()
// MAL:
props.initialValue = 50

// BIEN: copia local
const localValue = ref(props.initialValue)
\end{lstlisting}
\end{trampa}

% ══════════════════════════════════════════════════════════════════════════════
% NIVEL 4: COMPOSABLES
% ══════════════════════════════════════════════════════════════════════════════
\section{Nivel 4 — Composables: Pergaminos Reutilizables}

\begin{lore}
Los composables son la magia que todo Mago de Vue aprende en el Nivel 4.
Un composable encapsula lógica stateful reutilizable: se lleva el estado,
el ciclo de vida y las funciones juntos en un pergamino que cualquier
componente puede invocar.
\end{lore}

\subsection{Anatomía de un composable}

Un composable es simplemente una función que empieza con \texttt{use}
y utiliza la API de composición de Vue internamente.

\begin{pergamino}
\begin{lstlisting}[style=vue]
// composables/useCharges.ts
import { ref, computed, onMounted } from 'vue'

interface Charge {
  id: string
  amount: number
  status: string
}

export function useCharges(mandateId?: string) {
  const charges = ref<Charge[]>([])
  const isLoading = ref(false)
  const error = ref<string | null>(null)

  // Estado derivado
  const total = computed(() =>
    charges.value.reduce((sum, c) => sum + c.amount, 0)
  )
  const succeeded = computed(() =>
    charges.value.filter(c => c.status === 'succeeded')
  )
  const failed = computed(() =>
    charges.value.filter(c => c.status === 'failed')
  )

  async function load() {
    isLoading.value = true
    error.value = null
    try {
      const url = mandateId
        ? `/api/charges?mandate=${mandateId}`
        : '/api/charges'
      const resp = await fetch(url)
      if (!resp.ok) throw new Error(await resp.text())
      charges.value = await resp.json()
    } catch (e) {
      error.value = e instanceof Error
        ? e.message
        : 'Error desconocido'
    } finally {
      isLoading.value = false
    }
  }

  async function retry(chargeId: string) {
    await fetch(`/api/charges/${chargeId}/retry`, {
      method: 'POST'
    })
    await load()  // recargar lista
  }

  // Cargar automaticamente al montar
  onMounted(load)

  // Exponer lo que el componente necesita
  return {
    charges,
    isLoading,
    error,
    total,
    succeeded,
    failed,
    load,
    retry
  }
}
\end{lstlisting}
\end{pergamino}

\begin{pergamino}
\begin{lstlisting}[style=vue]
<!-- Uso en cualquier componente -->
<script setup lang="ts">
import { useCharges } from '@/composables/useCharges'

// Sin argumento: todos los cobros
const { charges, isLoading, error, total } = useCharges()

// Con argumento: cobros de un mandato
const {
  charges: mandateCharges,
  retry
} = useCharges('MAND-001')
</script>

<template>
  <div v-if="isLoading">Cargando...</div>
  <div v-else-if="error" class="error">{{ error }}</div>
  <ul v-else>
    <li v-for="charge in charges" :key="charge.id">
      {{ charge.id }} -- {{ charge.amount }}
      <button @click="retry(charge.id)">Reintentar</button>
    </li>
  </ul>
  <p>Total: {{ total }}</p>
</template>
\end{lstlisting}
\end{pergamino}

\subsection{Composable de formulario con validación}

\begin{pergamino}
\begin{lstlisting}[style=vue]
// composables/useForm.ts
import { reactive, ref, computed } from 'vue'

type Rules<T> = {
  [K in keyof T]?: (value: T[K]) => string | null
}

export function useForm<T extends object>(
  initialValues: T,
  rules?: Rules<T>
) {
  const values = reactive({ ...initialValues }) as T
  const errors = reactive({} as Record<keyof T, string | null>)
  const touched = reactive({} as Record<keyof T, boolean>)
  const isSubmitting = ref(false)

  function validate(): boolean {
    if (!rules) return true
    let valid = true
    for (const key in rules) {
      const rule = rules[key]
      if (rule) {
        const err = rule(values[key])
        errors[key] = err
        if (err) valid = false
      }
    }
    return valid
  }

  const isValid = computed(() =>
    !Object.values(errors).some(Boolean)
  )

  async function handleSubmit(
    onSubmit: (values: T) => Promise<void>
  ) {
    if (!validate()) return
    isSubmitting.value = true
    try {
      await onSubmit({ ...values })
    } finally {
      isSubmitting.value = false
    }
  }

  return { values, errors, touched,
           isValid, isSubmitting, validate,
           handleSubmit }
}
\end{lstlisting}
\end{pergamino}

% ══════════════════════════════════════════════════════════════════════════════
% NIVEL 5: PINIA
% ══════════════════════════════════════════════════════════════════════════════
\section{Nivel 5 — Pinia: el Grimorio del Estado Global}

\begin{lore}
Pinia es el grimorio oficial de Vue 3. Cuando el estado necesita
ser compartido entre componentes que no tienen relación directa
(no padre-hijo), o cuando persiste entre rutas, es hora de
invocar a Pinia. Reemplazó a Vuex como el grimorio estándar.
\end{lore}

\subsection{Definir un store}

\begin{pergamino}
\begin{lstlisting}[style=vue]
// stores/useChargeStore.ts
import { defineStore } from 'pinia'
import { ref, computed } from 'vue'

// Setup store (formato moderno, igual a composable)
export const useChargeStore = defineStore('charges', () => {
  // State
  const charges = ref<Charge[]>([])
  const isLoading = ref(false)
  const selectedId = ref<string | null>(null)

  // Getters (computed)
  const selected = computed(() =>
    charges.value.find(c => c.id === selectedId.value)
  )
  const total = computed(() =>
    charges.value.reduce((sum, c) => sum + c.amount, 0)
  )
  const pendingCount = computed(() =>
    charges.value.filter(c => c.status === 'pending').length
  )

  // Actions
  async function fetchAll() {
    isLoading.value = true
    try {
      const resp = await fetch('/api/charges')
      charges.value = await resp.json()
    } finally {
      isLoading.value = false
    }
  }

  async function retryCharge(id: string) {
    await fetch(`/api/charges/${id}/retry`, {
      method: 'POST'
    })
    // Actualizar localmente (optimistic update)
    const charge = charges.value.find(c => c.id === id)
    if (charge) charge.status = 'processing'
    // Recargar desde el servidor
    await fetchAll()
  }

  function select(id: string) {
    selectedId.value = id
  }

  return {
    charges, isLoading, selectedId,
    selected, total, pendingCount,
    fetchAll, retryCharge, select
  }
})

// Para Hot Module Replacement en desarrollo
if (import.meta.hot) {
  import.meta.hot.accept(
    acceptHMRUpdate(useChargeStore, import.meta.hot)
  )
}
\end{lstlisting}
\end{pergamino}

\begin{pergamino}
\begin{lstlisting}[style=vue]
<!-- Uso en componentes -->
<script setup lang="ts">
import { useChargeStore } from '@/stores/useChargeStore'
import { storeToRefs } from 'pinia'

const store = useChargeStore()

// storeToRefs: desestructurar sin perder reactividad
// (solo para state y getters, NO para actions)
const { charges, isLoading, total, pendingCount }
  = storeToRefs(store)

// Actions: desestructurar directamente
const { fetchAll, retryCharge, select } = store

// Llamar al montar
onMounted(fetchAll)
</script>

<template>
  <p>{{ pendingCount }} cobros pendientes</p>
  <p>Total: {{ total }}</p>
  <div v-if="isLoading">Cargando...</div>
  <ul v-else>
    <li v-for="charge in charges" :key="charge.id">
      <button @click="select(charge.id)">
        {{ charge.id }}
      </button>
      <button @click="retryCharge(charge.id)">
        Reintentar
      </button>
    </li>
  </ul>
</template>
\end{lstlisting}
\end{pergamino}

\begin{cofre}
\textbf{¿Composable o Pinia?}
Usa un \textbf{composable} cuando el estado es local a un árbol de
componentes o a una ruta. Usa \textbf{Pinia} cuando el estado necesita
persistir entre rutas, ser compartido entre árboles no relacionados
o ser accedido desde fuera de componentes (ej. interceptors de axios).
\end{cofre}

% ══════════════════════════════════════════════════════════════════════════════
% NIVEL 6: SLOTS Y PATRONES AVANZADOS
% ══════════════════════════════════════════════════════════════════════════════
\section{Nivel 6 — Slots y Patrones Avanzados}

\subsection{Slots — Las Puertas del Héroe}

Los slots permiten que el padre inyecte contenido dentro del hijo,
creando componentes flexibles y reutilizables.

\begin{pergamino}
\begin{lstlisting}[style=vue]
<!-- BaseCard.vue: componente generico con slots -->
<script setup lang="ts">
defineProps<{
  title: string
  variant?: 'default' | 'danger' | 'success'
}>()
</script>

<template>
  <div class="card" :class="`card--${variant ?? 'default'}`">
    <!-- Slot nombrado para el header -->
    <header>
      <h3>{{ title }}</h3>
      <slot name="actions" />
    </header>

    <!-- Slot default: contenido principal -->
    <main>
      <slot />
    </main>

    <!-- Slot opcional con fallback -->
    <footer>
      <slot name="footer">
        <p class="muted">Sin acciones adicionales</p>
      </slot>
    </footer>
  </div>
</template>
\end{lstlisting}
\end{pergamino}

\begin{pergamino}
\begin{lstlisting}[style=vue]
<!-- Uso de BaseCard -->
<template>
  <BaseCard title="Cobro #123" variant="success">
    <!-- Slot nombrado: template con #nombre -->
    <template #actions>
      <button @click="exportPDF()">PDF</button>
      <button @click="retry()">Reintentar</button>
    </template>

    <!-- Slot default: contenido directo -->
    <ChargeDetails :charge="charge" />

    <!-- Slot footer -->
    <template #footer>
      <span>Procesado el {{ charge.createdAt }}</span>
    </template>
  </BaseCard>
</template>
\end{lstlisting}
\end{pergamino}

\subsection{Scoped Slots — El Pergamino que Habla de Vuelta}

\begin{pergamino}
\begin{lstlisting}[style=vue]
<!-- DataTable.vue: lista generica con scoped slot -->
<script setup lang="ts" generic="T">
defineProps<{ items: T[] }>()
</script>

<template>
  <table>
    <tbody>
      <tr v-for="(item, index) in items" :key="index">
        <!-- Scoped slot: expone item y index al padre -->
        <slot :item="item" :index="index" />
      </tr>
    </tbody>
  </table>
</template>
\end{lstlisting}
\end{pergamino}

\begin{pergamino}
\begin{lstlisting}[style=vue]
<!-- Uso: el padre decide como renderizar cada fila -->
<DataTable :items="charges">
  <template #default="{ item: charge, index }">
    <td>{{ index + 1 }}</td>
    <td>{{ charge.id }}</td>
    <td>{{ charge.amount }}</td>
    <td>
      <span :class="charge.status">
        {{ charge.status }}
      </span>
    </td>
  </template>
</DataTable>
\end{lstlisting}
\end{pergamino}

\subsection{provide / inject — El Oráculo de la Torre}

\begin{pergamino}
\begin{lstlisting}[style=vue]
// Patron: proveer contexto desde un componente raiz
// sin prop drilling (pasar props nivel a nivel)

// AppRoot.vue o el componente raiz del modulo
<script setup lang="ts">
import { provide, ref } from 'vue'

const currentUser = ref<User | null>(null)
const theme = ref<'light' | 'dark'>('light')

// Clave tipada para evitar colisiones
provide('currentUser', currentUser)
provide('theme', theme)
</script>

// --- Cualquier descendiente, a cualquier nivel ---
<script setup lang="ts">
import { inject, ref, Ref } from 'vue'

// inject con tipo y valor por defecto
const currentUser = inject<Ref<User | null>>(
  'currentUser',
  ref(null)
)

const theme = inject<Ref<'light' | 'dark'>>(
  'theme',
  ref('light')
)
</script>
\end{lstlisting}
\end{pergamino}

\begin{cofre}
\textbf{provide/inject vs Pinia:} Usa \texttt{provide/inject} para
datos de contexto del componente (tema, usuario activo en un módulo).
Usa Pinia para estado global de aplicación que persiste entre rutas.
\end{cofre}

\subsection{Directivas esenciales en el template}

\begin{center}
\begin{tabular}{p{3.5cm} p{9cm}}
\toprule
\textbf{Directiva} & \textbf{Uso} \\
\midrule
\texttt{v-if / v-else-if / v-else}
  & Renderizado condicional (elimina el DOM) \\
\texttt{v-show}
  & Ocultar con \texttt{display:none} (mantiene el DOM) \\
\texttt{v-for="x in xs" :key}
  & Renderizado de listas. \texttt{:key} es obligatorio \\
\texttt{v-model}
  & Two-way binding en inputs, selects, checkboxes \\
\texttt{:prop="valor"}
  & Bind dinámico de atributos/props \\
\texttt{@evento="handler"}
  & Escuchar eventos del DOM o emits de hijo \\
\texttt{@click.prevent}
  & Modificadores: prevent, stop, once, self \\
\texttt{v-memo="[dep]"}
  & Memorización de nodo (optimización) \\
\bottomrule
\end{tabular}
\end{center}

% ══════════════════════════════════════════════════════════════════════════════
% NIVEL 7: MISIONES
% ══════════════════════════════════════════════════════════════════════════════
\section{Nivel 7 — Misiones Finales}

\begin{mision}{1 — La Joya del Formulario}
Crea un componente \texttt{MandateForm.vue} que permita crear un mandato
bancario. Debe tener campos para email, número de cuenta y monto máximo.
El botón de envío se deshabilita si algún campo está vacío o el monto
es inválido. Al enviar exitosamente, emite \texttt{'mandate-created'} con
los datos.
\end{mision}

\begin{recompensa}{1}
\begin{lstlisting}[style=vue]
<script setup lang="ts">
import { reactive, computed } from 'vue'

const emit = defineEmits<{
  'mandate-created': [mandate: Mandate]
}>()

const form = reactive({
  email: '', accountNumber: '', maxAmount: ''
})
const isSubmitting = ref(false)

const isValid = computed(() =>
  form.email.includes('@') &&
  form.accountNumber.length >= 6 &&
  Number(form.maxAmount) > 0
)

async function submit() {
  if (!isValid.value) return
  isSubmitting.value = true
  try {
    const resp = await fetch('/api/mandates', {
      method: 'POST',
      headers: { 'Content-Type': 'application/json' },
      body: JSON.stringify({
        email: form.email,
        accountNumber: form.accountNumber,
        maxAmount: Math.round(
          Number(form.maxAmount) * 100
        )
      })
    })
    const mandate = await resp.json()
    emit('mandate-created', mandate)
  } finally {
    isSubmitting.value = false
  }
}
</script>

<template>
  <form @submit.prevent="submit">
    <input v-model="form.email"
           type="email" placeholder="Email" />
    <input v-model="form.accountNumber"
           placeholder="N. de cuenta" />
    <input v-model="form.maxAmount"
           type="number" placeholder="Monto max ($)" />
    <button :disabled="!isValid || isSubmitting">
      {{ isSubmitting ? 'Creando...' : 'Crear mandato' }}
    </button>
  </form>
</template>
\end{lstlisting}
\end{recompensa}

\begin{mision}{2 — El Composable Encantado}
Crea un composable \texttt{usePagination<T>} que reciba una URL de API
y exponga: \texttt{items}, \texttt{page}, \texttt{totalPages},
\texttt{isLoading}, \texttt{nextPage()}, \texttt{prevPage()}.
Debe cargar automáticamente al cambiar de página.
\end{mision}

\begin{recompensa}{2}
\begin{lstlisting}[style=vue]
// composables/usePagination.ts
import { ref, computed, watch } from 'vue'

export function usePagination<T>(
  baseUrl: string,
  pageSize: number = 10
) {
  const items = ref<T[]>([])
  const page = ref(1)
  const total = ref(0)
  const isLoading = ref(false)

  const totalPages = computed(() =>
    Math.ceil(total.value / pageSize)
  )
  const hasPrev = computed(() => page.value > 1)
  const hasNext = computed(
    () => page.value < totalPages.value
  )

  async function load() {
    isLoading.value = true
    try {
      const url =
        `${baseUrl}?page=${page.value}` +
        `&limit=${pageSize}`
      const resp = await fetch(url)
      const data = await resp.json()
      items.value = data.items
      total.value = data.total
    } finally {
      isLoading.value = false
    }
  }

  // Recargar al cambiar de pagina
  watch(page, load, { immediate: true })

  const nextPage = () => {
    if (hasNext.value) page.value++
  }
  const prevPage = () => {
    if (hasPrev.value) page.value--
  }

  return {
    items, page, total, totalPages,
    isLoading, hasPrev, hasNext,
    nextPage, prevPage
  }
}
\end{lstlisting}
\end{recompensa}

\begin{mision}{3 — El Grimorio en Acción}
Modifica el store \texttt{useChargeStore} para que soporte
\textbf{optimistic updates}: al hacer retry, actualiza el estado
local inmediatamente a \texttt{'processing'} y revierte a
\texttt{'failed'} si la API falla.
\end{mision}

\begin{recompensa}{3}
\begin{lstlisting}[style=vue]
// En useChargeStore.ts -- action retryCharge mejorada
async function retryCharge(id: string) {
  const charge = charges.value.find(c => c.id === id)
  if (!charge) return

  // Guardar estado original
  const originalStatus = charge.status

  // Optimistic update: actualizar UI inmediatamente
  charge.status = 'processing'

  try {
    await fetch(`/api/charges/${id}/retry`, {
      method: 'POST'
    })
    // Recargar estado real del servidor
    await fetchAll()
  } catch {
    // Revertir si falla
    charge.status = originalStatus
    throw new Error('No se pudo reintentar el cobro')
  }
}
\end{lstlisting}
\end{recompensa}

% ══════════════════════════════════════════════════════════════════════════════
% CHEATSHEET
% ══════════════════════════════════════════════════════════════════════════════
\newpage
\section{El Pergamino Mayor — Cheatsheet Vue 3}

\begin{multicols}{2}

\textbf{\textcolor{vuegreen}{Reactividad}}
\begin{lstlisting}[style=vue,numbers=none]
const x = ref<number>(0)
x.value++          // modificar
const { a, b } = toRefs(reactive({a:1,b:2}))
const c = computed(() => x.value * 2)
\end{lstlisting}

\textbf{\textcolor{vuegreen}{Ciclo de vida}}
\begin{lstlisting}[style=vue,numbers=none]
onMounted(async () => { await load() })
onUnmounted(() => { clearInterval(timer) })
\end{lstlisting}

\textbf{\textcolor{vuegreen}{Watchers}}
\begin{lstlisting}[style=vue,numbers=none]
watch(x, (nuevo, viejo) => { ... })
watch([x, y], ([nx, ny]) => { ... })
watchEffect(() => { /* usa x.value */ })
\end{lstlisting}

\textbf{\textcolor{vuegreen}{Componente}}
\begin{lstlisting}[style=vue,numbers=none]
const props = defineProps<{ id: string }>()
const emit = defineEmits<{
  'event': [payload: string]
}>()
defineExpose({ publicMethod })
\end{lstlisting}

\columnbreak

\textbf{\textcolor{golddnd}{Pinia}}
\begin{lstlisting}[style=vue,numbers=none]
const store = defineStore('name', () => {
  const x = ref(0)         // state
  const y = computed(...)  // getter
  async function act() {}  // action
  return { x, y, act }
})
// En componente:
const s = useMyStore()
const { x } = storeToRefs(s)
s.act()
\end{lstlisting}

\textbf{\textcolor{vuegreen}{Template}}
\begin{lstlisting}[style=vue,numbers=none]
v-if="cond"
v-for="x in xs" :key="x.id"
v-model="ref"
:prop="expr"
@click="fn"
@click.prevent.stop
<slot name="header" />
<template #header>...</template>
\end{lstlisting}

\textbf{\textcolor{purplemage}{Composable pattern}}
\begin{lstlisting}[style=vue,numbers=none]
export function useX() {
  const state = ref(...)
  onMounted(load)
  async function load() { ... }
  return { state, load }
}
\end{lstlisting}

\end{multicols}

\vspace{1em}
\noindent\textcolor{vuegreen}{\rule{\linewidth}{0.8pt}}
\vspace{0.5em}

\begin{center}
\begin{tabular}{p{5cm} p{8cm}}
\toprule
\textbf{Pregunta} & \textbf{Respuesta clave} \\
\midrule
\texttt{ref} vs \texttt{reactive}
  & \texttt{ref} para primitivos y pasar referencias;
    \texttt{reactive} para objetos de estado agrupado \\
\texttt{computed} vs \texttt{watch}
  & \texttt{computed} deriva un valor (sin efectos);
    \texttt{watch} reacciona a cambios con efectos \\
Composable vs Pinia
  & Composable = estado local al árbol;
    Pinia = estado global entre rutas \\
\texttt{v-if} vs \texttt{v-show}
  & \texttt{v-if} elimina el DOM;
    \texttt{v-show} oculta con CSS \\
¿Por qué \texttt{:key} en \texttt{v-for}?
  & Vue lo usa para identificar nodos y hacer
    reconciliación eficiente del DOM \\
\bottomrule
\end{tabular}
\end{center}

\end{document}
